% Benchmark calculation paragraph for results.tex

% Formal version

To construct a benchmark for redistribution, we use aggregate data from the Nilson Report on interchange fees and payment volumes.\footnote{The Nilson Report reports average merchant fees of 2.25\% for credit cards and 0.72\% for debit cards.
We estimate interchange fees by subtracting approximately 25 basis points for network fees and acquirer margins, yielding effective interchange rates of 2\% for credit and 0.4\% for debit.
Applied to transaction volumes of approximately \$4 trillion each for credit and debit cards, this implies \$80 billion in credit card interchange fees and \$16 billion in debit card interchange fees.} In 2022, we estimate U.S. merchants paid approximately \$80 billion in credit card interchange fees.
With credit cards representing roughly half of consumer payment volume, the remaining half---cash and debit users---collectively pay for half of credit card interchange fees through higher prices, while receiving none of the associated rewards.
This implies a net transfer of approximately \$40 billion annually from cash and debit users to credit card users.
This approach mirrors the methodology of \citet{SchuhShyStavins2010}, who estimated per-household transfers but did not report aggregate figures.

% More colloquial version

Public debate regarding interchange fees is currently polarized between two extreme views that obscure the underlying economic mechanism.
Merchant groups and policymakers frequently characterize the fees as a pure ``deadweight loss'' or ``hidden tax,'' often citing the total aggregate cost---over \$100 billion annually---as money siphoned out of the economy.
Conversely, banking industry advocates frame interchange fees solely as the necessary funding source for popular consumer rewards, ignoring the fact that these rewards are funded by higher retail prices paid by all consumers.
Both perspectives fail to address the central economic reality: interchange fees function as a regressive transfer system, not a pure cost or a free benefit.

To quantify this mechanism, we first construct a standard benchmark using aggregate data following the logic of previous literature.
Using 2022 data from the Nilson Report, we estimate U.S. merchants paid approximately \$80 billion in credit card interchange fees.
If we assume these costs are passed through uniformly to a consumer base where credit cards represent roughly half of payment volume, cash and debit users cross-subsidize credit users by approximately \$40 billion annually.
However, even this benchmark relies on the strong assumption that card and cash users shop at the same stores and pay the same prices.
By incorporating novel data on consumer sorting and merchant fee heterogeneity, our findings refine this estimate to \$30 billion.
This demonstrates that while market features like sector discounts reduce the transfer, they do not eliminate the substantial regressive cross-subsidy that defines the U.S. payments system.
